%%%%%%%%%%%%%%%%%%%%%%%%%%%%%%%%%%%%%%%%%%%%
% En 'inclues.tex' se encuentran la importación de paquetes necesarios
%%%%%%%%%%%%%%%%%%%%%%%%%%%%%%%%%%%%%%%%%%%%
%%%%%%%%%%%%%%%%%%%%%%%%%%%%%%%%%%%%%%%%%
% University Assignment Title Page 
% LaTeX Template
% Version 1.0 (27/12/12)
%
% This template has been downloaded from:
% http://www.LaTeXTemplates.com
%
% Original author:
% WikiBooks (http://en.wikibooks.org/wiki/LaTeX/Title_Creation)
%
% License:
% CC BY-NC-SA 3.0 (http://creativecommons.org/licenses/by-nc-sa/3.0/)
% 
% Instructions for using this template:
% This title page is capable of being compiled as is. This is not useful for 
% including it in another document. To do this, you have two options: 
%
% 1) Copy/paste everything between \begin{document} and \end{document} 
% starting at \begin{titlepage} and paste this into another LaTeX file where you 
% want your title page.
% OR
% 2) Remove everything outside the \begin{titlepage} and \end{titlepage} and 
% move this file to the same directory as the LaTeX file you wish to add it to. 
% Then add \input{./title_page_1.tex} to your LaTeX file where you want your
% title page.
%
%%%%%%%%%%%%%%%%%%%%%%%%%%%%%%%%%%%%%%%%%
%\title{Title page with logo}
%----------------------------------------------------------------------------------------
%	PACKAGES AND OTHER DOCUMENT CONFIGURATIONS
%----------------------------------------------------------------------------------------
\documentclass[14pt]{extarticle}
%Paquetes para idioma español y codifcación UTF8
\usepackage[spanish]{babel}
\usepackage[utf8x]{inputenc}
%natbib, para bibliografía y referencias
\usepackage[numbers]{natbib}
%paquete de matemática, vuela
\usepackage{amsmath}
%fuente 'fourier'
\usepackage{fourier}
%paquete para URLs
\usepackage{url}
%paquete para ubicar las imágenes
\usepackage{float}
%paquete para imágenes y en dónde las tiene que buscar
\usepackage{graphicx}
\graphicspath{{images/}}
%paquete para epígrafes
\usepackage{subcaption}
%paquete para definir los márgenes de la hoja
\usepackage[left=2cm,right=2cm,top=4cm,bottom=4cm]{geometry}
%paquete para poner todos y comentarios
\usepackage[colorinlistoftodos]{todonotes}
%paquete para trabajar con código
\usepackage{listings}
%paquete para trabajar con colores y definir propios
\usepackage{color}

%Cabeceras
\usepackage{fancyhdr}
\pagestyle{fancy}
\fancyhead[L]{Bases de Datos II, 2018}
\fancyhead[C]{}
\fancyhead[R]{UNPSJB}

\fancyfoot[R]{Luciano Serruya Aloisi}
\fancyfoot[L]{Trabajo Práctico 2}

%%paquete para reducir espacio entre secciones
%\usepackage{titlesec}
%\titlespacing*{\section}{0pt}{1.1\baselineskip}{\baselineskip}

%Comando para poner doble comillas más fácil
\newcommand{\dq}[1]{``#1''}

\definecolor{mygreen}{rgb}{0,0.6,0}
\definecolor{mygray}{rgb}{0.5,0.5,0.5}
\definecolor{mymauve}{rgb}{0.58,0,0.82}
\definecolor{dkgreen}{rgb}{0,0.6,0}
\definecolor{ltgray}{rgb}{0.5,0.5,0.5}
\definecolor{alizarin}{rgb}{0.82, 0.1, 0.26}

\lstdefinestyle{customstyle}{ 
  backgroundcolor=\color{white},   % choose the background color; you must add \usepackage{color} or \usepackage{xcolor}; should come as last argument
  basicstyle=\footnotesize,        % the size of the fonts that are used for the code
  breakatwhitespace=false,         % sets if automatic breaks should only happen at whitespace
  breaklines=true,                 % sets automatic line breaking
  captionpos=b,                    % sets the caption-position to bottom
  commentstyle=\color{mygreen},    % comment style
  escapeinside={\%*}{*)},          % if you want to add LaTeX within your code
  extendedchars=true,              % lets you use non-ASCII characters; for 8-bits encodings only, does not work with UTF-8
  frame=single,	                   % adds a frame around the code
  keepspaces=true,                 % keeps spaces in text, useful for keeping indentation of code (possibly needs columns=flexible)
  keywordstyle=\color{blue},       % keyword style
  language=Octave,                 % the language of the code
  numbers=left,                    % where to put the line-numbers; possible values are (none, left, right)
  numbersep=5pt,                   % how far the line-numbers are from the code
  numberstyle=\tiny\color{mygray}, % the style that is used for the line-numbers
  rulecolor=\color{black},         % if not set, the frame-color may be changed on line-breaks within not-black text (e.g. comments (green here))
  showspaces=false,                % show spaces everywhere adding particular underscores; it overrides 'showstringspaces'
  showstringspaces=false,          % underline spaces within strings only
  showtabs=false,                  % show tabs within strings adding particular underscores
  stepnumber=2,                    % the step between two line-numbers. If it's 1, each line will be numbered
  stringstyle=\color{mymauve},     % string literal style
  tabsize=2,	                   % sets default tabsize to 2 spaces
  title=\lstname                   % show the filename of files included with \lstinputlisting; also try caption instead of title
}

\lstdefinestyle{sql}{
	basicstyle=\fontsize{12}{10}\ttfamily,
	belowcaptionskip=1\baselineskip,
  	breaklines=true,
  	title=\lstname,	% show the filename of files included with \lstinputlisting; also try caption instead of title
  	frame=L, %Este pone la línea (doble) separadora a la izquiera, ir probando entre l, L,r y R
  	xleftmargin=\parindent,
  	language=SQL,
  	morekeywords={*,modify,MODIFY,...},
  	numbers=left,                    % where to put the line-numbers; possible values are (none, left, right)
  	numberstyle=\color{mygray}, % the style that is used for the line-numbers
  	rulecolor=\color{mygray},         % if not set, the frame-color may be changed on line-breaks within not-black text (e.g. comments (green here))
  	showstringspaces=false,
  	captionpos=b,                    % sets the caption-position to bottom
  	keywordstyle=\bfseries\color{green!40!black},
  	commentstyle=\itshape\color{purple!40!black},
  	identifierstyle=\color{blue},
  	stringstyle=\color{alizarin},
  	escapechar=@,
}

\lstdefinestyle{cli}{
    basicstyle=\fontsize{12}{10}\ttfamily,%\fontfamily{phv}\selectfont,
	belowcaptionskip=1\baselineskip,
  	breaklines=true,
  	title=\lstname,	% show the filename of files included with \lstinputlisting; also try caption instead of title
  	%frame=L, %Este pone la línea (doble) separadora a la izquiera, ir probando entre l, L,r y R
  	xleftmargin=\parindent,
  	language=bash,
  	%numbers=left,                    % where to put the line-numbers; possible values are (none, left, right)
  	%numberstyle=\color{mygray}, % the style that is used for the line-numbers
  	%rulecolor=\color{mygray},         % if not set, the frame-color may be changed on line-breaks within not-black text (e.g. comments (green here))
  	showstringspaces=false,
  	captionpos=b,                    % sets the caption-position to bottom
  	keywordstyle=\bfseries\color{blue!40!black},
  	commentstyle=\itshape\color{purple!40!black},
  	identifierstyle=\color{blue},
  	stringstyle=\color{alizarin},
  	escapechar=@,
}

\lstdefinestyle{bash}{
	basicstyle=\fontsize{12}{10}\ttfamily,
	belowcaptionskip=1\baselineskip,
  	breaklines=true,
  	%title=\lstname,	% show the filename of files included with \lstinputlisting; also try caption instead of title
  	%frame=L, %Este pone la línea (doble) separadora a la izquiera, ir probando entre l, L,r y R
  	xleftmargin=\parindent,
  	language=bash,
  	%numbers=left,                    % where to put the line-numbers; possible values are (none, left, right)
  	numberstyle=\color{mygray}, % the style that is used for the line-numbers
  	rulecolor=\color{mygray},         % if not set, the frame-color may be changed on line-breaks within not-black text (e.g. comments (green here))
  	showstringspaces=false,
  	captionpos=b,                    % sets the caption-position to bottom
  	keywordstyle=\bfseries\color{green!40!black},
  	commentstyle=\itshape\color{purple!40!black},
  	identifierstyle=\color{blue},
  	stringstyle=\color{alizarin},
  	escapechar=@,
}



\begin{document}

%%%%%%%%%%%%%%%%%%%%%%%%%%%%%%%%%%%%%%%%%%%%
% En 'titlepage.tex' se encuentra la página de título
%%%%%%%%%%%%%%%%%%%%%%%%%%%%%%%%%%%%%%%%%%%%
\begin{titlepage}

    \newcommand{\HRule}{\rule{\linewidth}{0.5mm}} % Defines a new command for the horizontal lines, change thickness here

    \center % Center everything on the page
     
    %----------------------------------------------------------------------------------------
    %	HEADING SECTIONS
    %----------------------------------------------------------------------------------------

    \textsc{\LARGE UNPSJB}\\[1cm] % Name of your university/college
    \textsc{\Large Licenciatura en Sistemas OPGCPI}\\[0.5cm] % Major heading such as course name
    \textsc{\large Bases de Datos II}\\[0.5cm] % Minor heading such as course title

    %----------------------------------------------------------------------------------------
    %	TITLE SECTION
    %----------------------------------------------------------------------------------------

    \HRule \\[0.4cm]
    {\huge \bfseries Trabajo Práctico 2}\\[0.4cm] % Title of your document
    {\large \bfseries Postgres, PL/pgSQL, Integridad y Seguridad de los Datos}\\[0.4cm] % Title of your document
    \HRule \\[1.5cm]
     
    %----------------------------------------------------------------------------------------
    %	AUTHOR SECTION
    %----------------------------------------------------------------------------------------


    \begin{minipage}[l]{0.4\textwidth}
        \begin{flushleft}
            \textbf{\textsf{Cátedra}}\\
            \large Lic. Gabriel Ingravallo\\ 
            \large Lic. Cristian Parisse\\ 
            \linespread{4}
            \end{flushleft}
    \end{minipage}
    \begin{minipage}[l]{0.4\textwidth}
        \begin{flushright}
            \textbf{\textsf{Integrantes:}}\\
            \linespread{1}
            \large Luciano Serruya Aloisi
        \end{flushright}
    \end{minipage}\\[2cm]

    % If you don't want a supervisor, uncomment the two lines below and remove the section above
    %\Large \emph{Author:}\\
    %John \textsc{Smith}\\[3cm] % Your name

    %----------------------------------------------------------------------------------------
    %	DATE SECTION
    %----------------------------------------------------------------------------------------

    {\large \today}\\[1cm] % Date, change the \today to a set date if you want to be precise

    %----------------------------------------------------------------------------------------
    %	LOGO SECTION
    %----------------------------------------------------------------------------------------

    \includegraphics[scale=1]{logoUnpsjb.png}\\[0.5cm] % Include a department/university logo - this will require the graphicx package
     
    %----------------------------------------------------------------------------------------

    % \vfill % Fill the rest of the page with whitespace

\end{titlepage}


%%%%%%%%%%%%%%%%%%%%%%%%%%%%%%%%%%%%%%%%%%%%
% INDICE
%%%%%%%%%%%%%%%%%%%%%%%%%%%%%%%%%%%%%%%%%%%%
\clearpage
\tableofcontents
\clearpage

%%%%%%%%%%%%%%%%%%%%%%%%%%%%%%%%%%%%%%%%%%%%
% ABSTRACT
%%%%%%%%%%%%%%%%%%%%%%%%%%%%%%%%%%%%%%%%%%%%
%\begin{abstract}
%Your abstract.
%\end{abstract}

\section{Carga masiva de la Base de Datos}

\emph{Agregar tuplas en forma masiva a las tablas modeloAvion y Avion de forma tal que por cada modeloAvion insertado se agreguen 10000 aviones y que la estadística V(capacidad, modeloAvion) >= 1000. Además por cada avión agregado agregar una reparación del mismo (con algún DNI de trabajador ya preexistente). Se debe presentar el script de la solucion.} 

~

\emph{\textbf{Los scripts a los que se haga referencia se pueden encontrar en el directorio scripts/}} 

\subsection{Modelos de avión}

Los modelos de avión utilizados para la carga masiva de la base de datos fueron tomados del \emph{dataset} disponible en \cite{dataset:2013}. Los tipos de naves del \emph{dataset} fueron extraídos mediante el \emph{script} de Python \textbf{extraer\_modelos.py}

A continuación, mediante otro \emph{script} en Python (\textbf{generar\_script.py}), generamos un \emph{script} en SQL para hacer la carga de los modelos de avión a la base de datos (\textbf{script\_carga\_modelos.sql}). De esta forma, insertamos en la base el total de \textbf{2449} modelos de avión, siendo su clave primaria un número secuencial y su capacidad un valor aleatorio entre 50 y 2000. 

\vspace*{5mm}
\lstset{style=sql}
\begin{lstlisting}
aviones=# SELECT COUNT(*) FROM "modeloAvion";
 count 
-------
  2449
(1 row)
\end{lstlisting}

Los valores cargados para los modelos de avión cumplen con la estadística V(capacidad, modeloAvion) $\geq$ 1000

\vspace*{5mm}
\lstset{style=sql}
\begin{lstlisting}
aviones=# SELECT COUNT(DISTINCT capacidad) FROM "modeloAvion";
 count 
-------
  1389
(1 row)
\end{lstlisting}

\subsection{Trabajadores}

Para la carga de los trabajadores, se utilizó el sitio \cite{generatedata} utilizando valores aleatorios para los DNIs y nombres de ascendencia latina disponibles en el repositorio de \emph{GitHub} \cite{apellidos:2012}. El script generado por el sitio se encuentra bajo el nombre de \textbf{script\_carga\_trabajadores.sql} 

\vspace*{5mm}
\lstset{style=sql}
\begin{lstlisting}
aviones=# SELECT COUNT(*) FROM trabajador;
 count 
-------
   105
(1 row)
\end{lstlisting}


\subsection{Aviones}
Una vez ejecutados exitósamente los dos scripts de carga, se procede a insertar masivamente registros en la tabla de aviones. 

Por cada modelo de avión existente en la base de datos, se cargan 10000 aviones, con una cantidad de horas de vuelo aleatoria entre 100 y 5000, y un año aleatorio entre 1970 y 2018. El \emph{script} de carga correspondiente es \textbf{script\_carga\_aviones.sql}.

La función que realiza la carga de aviones está compuesto por dos bucles \emph{FOR .. LOOP} anidados; el más externo itera sobre los modelos de avión, mientras que el interno sobre enteros entre 1 y 10000. 

\vspace*{5mm}
\lstset{style=sql}
\begin{lstlisting}
aviones=# SELECT COUNT(*) FROM avion;
  count   
----------
 24490006
(1 row)
\end{lstlisting}

\subsection{Reparaciones}
Por último, se cargan las reparaciones de los aviones. El \emph{script} correspondiente (\textbf{script\_carga\_reparaciones.sql}) consiste también de dos bucles \emph{FOR..LOOP} anidados; el primero itera sobre todos los aviones existentes, mientras que el segundo itera sobre un rango de valores aleatorios entre 1 y 10, que corresponden a la cantidad de reparaciones registradas para ese avión (puede haber registrado entre 1 a 10 reparaciones). El DNI del trabajador también va a ser un valor aleatorio entre todos los registrados
\vspace*{5mm}
\lstset{style=sql}
\begin{lstlisting}
SELECT INTO dni_trabajador dni FROM trabajador ORDER BY random() LIMIT 1;
\end{lstlisting}

La fecha de inicio de la reparación se obtiene aleatoriamente entre un intervalo de tiempo, empezando desde 01/01/2010, hasta el 01/03/2018. Por último, la fecha de fin de reparación, que por más que no es obligatoria, se obtiene a partir de sumar una cantidad de días aleatoria a la fecha de inicio (entre 2 y 30 días).

Debido a los valores aleatorios, puede ocurrir que se intenten insertar dos veces una misma tupla (una reparación que tenga la misma clave primaria, compuesta del DNI del trabajador, el ID del avión, y la fecha de inicio de la reparación), por lo tanto se decidió que, frente a un intento de insertar un registro repetido, no hacer nada (no insertar la tupla conflictuante ni actualizar la existente).

\vspace*{5mm}
\lstset{style=sql}
\begin{lstlisting}
INSERT INTO "trabajadorReparacion" VALUES(dni_trabajador, nro_avion, fecha_inicio, fecha_fin, tipo_falla) ON CONFLICT DO NOTHING;
\end{lstlisting}

\lstset{style=sql}
\begin{lstlisting}
aviones=# SELECT COUNT(*) FROM "trabajadorReparacion";
   count   
-----------
 134712980
(1 row)
\end{lstlisting}







\clearpage
\bibliographystyle{bst/unified}
\bibliography{references}

\end{document}
\todo[inline, color=green!40]{This is an inline comment.}
